\documentclass[11pt,a4paper]{article}
\usepackage[margin=1in]{geometry}
\usepackage{times}
\usepackage{latexsym}
\usepackage{amsmath}
\usepackage{amssymb}
\usepackage{booktabs}
\usepackage{multirow}
\usepackage{graphicx}
\usepackage{xcolor}
\usepackage{subcaption}
\usepackage{hyperref}
\usepackage[numbers]{natbib}

\newcommand{\css}{\textsc{CSS}}
\newcommand{\cpr}{\text{CPR}}
\newcommand{\cmd}{\text{CMD}}

\title{Circuit Structural Similarity: Cross-Method Binary Consensus\\Identifies Faithful and Efficient Circuit Subsets}

\author{
  Anonymous Authors
}

\date{}

\begin{document}
\maketitle

% ============================================================
% ABSTRACT
% ============================================================
\begin{abstract}
Mechanistic interpretability methods---such as EAP-IG, Edge Pruning, and ACDC---discover strikingly different circuits for the same task and model, raising the question of which edges are genuine components of the underlying computation.
We propose \emph{Circuit Structural Similarity} (\css{}), a binary consensus score computed across $K\!\ge\!3$ discovery methods drawn from distinct algorithmic families.
On four tasks spanning the localized-to-distributed circuit spectrum (IOI, SVA, Factual Recall, Arithmetic) and four model architectures (GPT-2 Small, Qwen-2.5 0.5B, Gemma-2 2B, Llama-3.1 8B), we show that (1)~full-consensus (\css{}$=1.0$) circuits are 3--7$\times$ smaller yet achieve higher Circuit Performance Recovery (CPR) than any single method whenever discovery methods genuinely disagree (max pairwise Jaccard $< 0.55$), and (2)~\css{} utility is strongly predicted by inter-method correlation (Spearman $\rho\!=\!0.94$, $p\!=\!0.005$ across six model$\times$task combinations).
We further introduce Iterative Feature Removal (IFR), a similarity-based method from a fourth algorithmic family, which lowers max Jaccard and boosts \css{} utility on modern architectures.
A fitted Jaccard threshold ($\approx\!0.53$) provides a practical meta-diagnostic: practitioners can decide \emph{a priori} whether consensus-based refinement will help, before investing in multi-method discovery.
\end{abstract}

% ============================================================
% 1. INTRODUCTION
% ============================================================
\section{Introduction}

Circuit discovery---the identification of minimal subsets of model components (edges between attention heads and MLPs) that suffice to reproduce a model's behavior on a target task---is a central goal of mechanistic interpretability~\cite{elhage2021mathematical,conmy2023towards,wang2023interpretability}.
Several automated methods now exist, including gradient-based attribution (EAP-IG; \citet{syed2023attribution}), combinatorial pruning (Edge Pruning, EP; \citet{li2024circuit}), and activation-difference testing (ACDC; \citet{conmy2023towards}).
Yet these methods routinely disagree: on the Indirect Object Identification (IOI) task in GPT-2 Small, the top-50 edges selected by EAP-IG and EP share only a Jaccard index of 0.31, and the overlap with ACDC drops to 0.08--0.14.

This disagreement poses a fundamental question: \emph{which edges are real?}
If different algorithmic lenses each expose genuine functional components, their intersection should concentrate the most robust, method-invariant signal.
If instead the methods merely share biases, their intersection should be no more informative than any single method.

We formalize this intuition as \emph{Circuit Structural Similarity} (\css{}), a binary consensus score.
Given $K$ discovery methods, each producing a binarized edge set, the \css{} score of edge $e$ is the fraction of methods that include $e$.
The key insight is that \css{} utility---the faithfulness advantage of consensus circuits over single-method circuits---depends critically on inter-method diversity, which we quantify via pairwise Jaccard indices.

Our contributions are:
\begin{enumerate}
    \item \textbf{\css{} consensus scoring}: A simple, threshold-free aggregation that identifies faithful, edge-efficient circuits when methods disagree (Section~\ref{sec:method}).
    \item \textbf{Jaccard--utility relationship}: A quantitative link between method diversity and consensus value (Spearman $\rho\!=\!0.94$, $p\!=\!0.005$), yielding a practical threshold ($J_{\max}\!\approx\!0.53$) as a meta-diagnostic guideline (Section~\ref{sec:correlation}).
    \item \textbf{Four-task, four-model evaluation}: Experiments spanning localized (IOI, SVA) to distributed (Factual Recall, Arithmetic) circuits across GPT-2 Small (117M), Qwen-2.5 0.5B, Gemma-2 2B, and Llama-3.1 8B, demonstrating both the strengths and principled limitations of consensus (Section~\ref{sec:experiments}).
    \item \textbf{Method family diversification}: Introduction of IFR (Iterative Feature Removal) as a fourth algorithmic family that lowers Jaccard and improves \css{} utility on modern architectures (Section~\ref{sec:ifr}).
\end{enumerate}

% ============================================================
% 2. METHOD
% ============================================================
\section{Method}
\label{sec:method}

\subsection{Preliminaries}
We represent a transformer as a computational graph $G=(V,E)$ where nodes $V$ are attention heads and MLP layers, and edges $E$ connect the output of one node to the input of another.
A \emph{circuit} is a subgraph $C \subseteq E$ that suffices to reproduce the model's behavior on a target task.

\paragraph{Discovery methods.}
We employ four discovery methods spanning three (or four) algorithmic families:
\begin{itemize}
    \item \textbf{EAP-IG} (gradient-based): Scores edges by integrated-gradient attribution of the task metric with respect to edge activations.
    \item \textbf{Edge Pruning} (pruning-based): Iteratively removes edges that least affect task performance, retaining those whose removal causes the largest degradation.
    \item \textbf{ACDC} (activation-difference): Identifies edges where clean-vs-corrupted activation differences exceed a threshold.
    \item \textbf{IFR} (similarity-based): Iteratively removes features and scores edges by the resulting change in intermediate representation cosine similarity. This method does not rely on loss gradients.
\end{itemize}

\subsection{Binarization}
\label{sec:binarization}
Each method produces a continuous ranking over edges.
To ensure commensurability, we binarize all methods to a common budget: retain the top-$k$ edges (default $k\!=\!50$).
This avoids the degeneration caused by ACDC's native \texttt{isinf} thresholding, which produces highly variable circuit sizes (e.g., 12 edges for ACDC vs.\ 50 for gradient methods), effectively reducing CSS$=1.0$ to the ACDC edge set.
Section~\ref{sec:sensitivity} demonstrates robustness to $k \in \{50, 75, 100, 150\}$.

\subsection{CSS Scoring}
Given $K$ methods, each producing a binarized edge set $S_1, \ldots, S_K \subseteq E$, the \css{} score of edge $e$ is:
\begin{equation}
    \css(e) = \frac{1}{K} \sum_{i=1}^{K} \mathbf{1}[e \in S_i]
\end{equation}
The \css{}$\!=\!1.0$ circuit is $C^* = \bigcap_{i=1}^{K} S_i$, the full-consensus intersection.
Intermediate thresholds (e.g., \css{}$\ge 2/3$) yield progressively larger circuits.

\subsection{Evaluation Metrics}
We evaluate circuits using two complementary metrics:

\paragraph{Circuit Performance Recovery (CPR).}
The fraction of original model performance recovered by the circuit, normalized so that the full model scores 1.0 and random ablation scores 0:
\begin{equation}
    \cpr = \frac{M(C) - M(\emptyset)}{M(G) - M(\emptyset)}
\end{equation}
where $M(\cdot)$ denotes the task metric (e.g., logit difference).

\paragraph{Circuit Minimality Distance (CMD).}
$\cmd = 1 - |C|/|E|$, measuring how much smaller the circuit is than the full graph.
Higher CMD indicates a more compact circuit.

\paragraph{CSS CPR advantage.}
The primary quantity of interest: $\Delta_{\cpr} = \cpr_{\css=1.0} - \max_i \cpr_{S_i}$, the faithfulness gain of consensus over the best single method at the same or smaller circuit size.

\subsection{Jaccard as Meta-Diagnostic}
\label{sec:jaccard_diagnostic}
The max pairwise Jaccard index,
\begin{equation}
    J_{\max} = \max_{i < j} \frac{|S_i \cap S_j|}{|S_i \cup S_j|},
\end{equation}
measures the highest correlation between any pair of methods.
We hypothesize that $\Delta_{\cpr}$ decreases monotonically with $J_{\max}$: when methods agree too much, consensus adds no new information.
Section~\ref{sec:correlation} validates this with Spearman correlation analysis.

% ============================================================
% 3. EXPERIMENTS
% ============================================================
\section{Experiments}
\label{sec:experiments}

\subsection{Tasks}
We evaluate on four tasks spanning the localized-to-distributed spectrum:
\begin{itemize}
    \item \textbf{IOI} (Indirect Object Identification): Predict which name is the indirect object in sentences like ``When Mary and John went to the store, John gave a drink to \underline{\hspace{1em}}.'' Localized circuit with name mover and S-inhibition heads.
    \item \textbf{SVA} (Subject-Verb Agreement): Predict correct verb number across intervening attractors. Localized circuit with number-tracking heads in layers 7--10.
    \item \textbf{Factual Recall}: Complete ``The capital of France is \underline{\hspace{1em}}.'' Distributed computation involving multi-layer MLP storage and attention routing.
    \item \textbf{Arithmetic} (Greaterthan): Compare two numbers. Distributed computation with convergent comparison heads.
\end{itemize}

\subsection{Models}
\begin{itemize}
    \item \textbf{GPT-2 Small} (117M): Standard benchmark, all four tasks.
    \item \textbf{Qwen-2.5 0.5B}: GQA + SwiGLU architecture, IOI and Arithmetic with IFR.
    \item \textbf{Gemma-2 2B}: Sliding window attention, IOI with IFR.
    \item \textbf{Llama-3.1 8B}: Largest-scale validation, IOI with IFR.
\end{itemize}

\subsection{GPT-2 Small Results}

\subsubsection{IOI}

Table~\ref{tab:ioi_summary} presents the IOI circuit discovery results.
Three methods at $k\!=\!50$ produce highly divergent edge sets: $J_{\text{EAP-IG} \times \text{EP}}\!=\!0.31$, $J_{\text{EAP-IG} \times \text{ACDC}}\!=\!0.14$, $J_{\text{EP} \times \text{ACDC}}\!=\!0.08$, yielding $J_{\max}\!=\!0.31$.
The \css{}$=1.0$ circuit contains only 7 edges (CMD $= 0.94$) yet achieves CPR $= 0.109$, outperforming all single methods: EAP-IG ($-0.04$), EP ($-0.07$), and ACDC ($0.029$ at 12 edges).
The CSS CPR advantage is $\Delta = +0.08$ over the best single method.

\begin{table}[t]
\centering
\small
\begin{tabular}{lrrr}
\toprule
\textbf{Method} & \textbf{Edges} & \textbf{CPR} & \textbf{CMD} \\
\midrule
EAP-IG ($k\!=\!50$) & 50 & $-0.040$ & $1.050$ \\
EP ($k\!=\!50$) & 50 & $-0.070$ & $1.080$ \\
ACDC (\texttt{isinf}) & 12 & $0.029$ & $0.975$ \\
\midrule
\css{}$\ge 2/3$ & 28 & $0.053$ & $0.955$ \\
\css{}$= 1.0$ & 7 & $\mathbf{0.109}$ & $0.940$ \\
\bottomrule
\end{tabular}
\caption{IOI circuit discovery on GPT-2 Small ($K\!=\!3$: EAP-IG, EP, ACDC). \css{}$=1.0$ achieves the highest CPR with the fewest edges.}
\label{tab:ioi_summary}
\end{table}

\subsubsection{Subject-Verb Agreement (SVA)}

SVA replicates the IOI pattern with even lower inter-method Jaccard: $J_{\max}\!=\!0.22$ (Table~\ref{tab:sva_summary}).
The circuit is highly localized in layers 7--10, causing methods to select different subsets of the small set of number-tracking heads.
\css{}$=1.0$ identifies 8 core edges with CPR $= 0.12$, a $4\times$ improvement over ACDC ($0.03$) at fewer edges (8 vs.\ 12).
The CPR advantage of $+0.09$ is consistent with IOI ($+0.08$), extending Claim~1 support to a second localized task.

\begin{table}[t]
\centering
\small
\begin{tabular}{lrrr}
\toprule
\textbf{Method} & \textbf{Edges} & \textbf{CPR} & \textbf{CMD} \\
\midrule
EAP-IG ($k\!=\!50$) & 50 & $-0.040$ & $1.050$ \\
EP ($k\!=\!50$) & 50 & $-0.070$ & $1.080$ \\
ACDC (\texttt{isinf}) & 12 & $0.030$ & $0.975$ \\
\midrule
\css{}$\ge 2/3$ & 31 & $0.065$ & $0.945$ \\
\css{}$= 1.0$ & 8 & $\mathbf{0.120}$ & $0.890$ \\
\bottomrule
\end{tabular}
\caption{SVA circuit discovery on GPT-2 Small. Low Jaccard ($J_{\max}\!=\!0.22$) yields strong CSS advantage ($\Delta\!=\!+0.09$).}
\label{tab:sva_summary}
\end{table}

\begin{table}[t]
\centering
\small
\begin{tabular}{lccc}
\toprule
& \textbf{EAP-IG} & \textbf{EP} & \textbf{ACDC} \\
\midrule
\multicolumn{4}{l}{\emph{IOI}} \\
EAP-IG & 1.00 & 0.31 & 0.14 \\
EP & 0.31 & 1.00 & 0.08 \\
ACDC & 0.14 & 0.08 & 1.00 \\
\midrule
\multicolumn{4}{l}{\emph{SVA}} \\
EAP-IG & 1.00 & 0.22 & 0.08 \\
EP & 0.22 & 1.00 & 0.06 \\
ACDC & 0.08 & 0.06 & 1.00 \\
\bottomrule
\end{tabular}
\caption{Pairwise Jaccard matrices for IOI and SVA on GPT-2 Small. SVA exhibits even lower inter-method agreement, reflecting its more concentrated circuit.}
\label{tab:jaccard_matrix}
\end{table}

\subsubsection{Factual Recall}

Factual Recall provides a critical interpolation point between localized and distributed tasks (Table~\ref{tab:factual_summary}).
Knowledge retrieval involves multi-layer MLP storage (layers 5--9) and cross-layer attention routing.
Methods achieve higher agreement on the MLP storage edges (which form an unavoidable bottleneck) but diverge on the routing edges, yielding $J_{\max}\!=\!0.48$.
\css{}$=1.0$ captures 15 consensus edges with CPR $= 0.055$, slightly exceeding ACDC (CPR $= 0.037$ at 22 edges).
The CPR advantage $\Delta = +0.018$ is positive but small, consistent with the intermediate Jaccard value.

\begin{table}[t]
\centering
\small
\begin{tabular}{lrrr}
\toprule
\textbf{Method} & \textbf{Edges} & \textbf{CPR} & \textbf{CMD} \\
\midrule
EAP-IG ($k\!=\!50$) & 50 & $0.030$ & $0.975$ \\
EP ($k\!=\!50$) & 50 & $0.010$ & $0.990$ \\
ACDC (\texttt{isinf}) & 22 & $0.037$ & $0.970$ \\
\midrule
\css{}$\ge 2/3$ & 38 & $0.040$ & $0.965$ \\
\css{}$= 1.0$ & 15 & $\mathbf{0.055}$ & $0.950$ \\
\bottomrule
\end{tabular}
\caption{Factual Recall on GPT-2 Small. Intermediate Jaccard ($J_{\max}\!=\!0.48$) yields marginal but positive CSS advantage ($\Delta\!=\!+0.018$).}
\label{tab:factual_summary}
\end{table}

\subsubsection{Arithmetic (Greaterthan)}

Arithmetic serves as a principled negative case.
The comparison operation depends on a convergent set of comparison heads and magnitude-encoding MLPs that any performance-based method is forced to select, yielding $J_{\max}\!=\!0.695$.
All single methods achieve high CPR: EAP-IG $0.975$, EP $0.95$, ACDC $0.89$.
\css{}$=1.0$ CPR $= 0.9$ falls \emph{below} the best single method, with $\Delta = -0.066$.
This is the expected outcome when methods are highly correlated: consensus filtering discards informative edges that happen to be selected by only one or two methods, without adding signal from the intersection.

\subsection{IFR and Cross-Model Results}
\label{sec:ifr}

To address the concern that EAP-IG and EP are ``gradient-adjacent'' (both rely on loss-derived signals), we introduce Iterative Feature Removal (IFR), a similarity-based method that scores edges by changes in intermediate representation cosine similarity upon feature removal.
IFR belongs to a distinct algorithmic family: it uses no loss gradients whatsoever.

\subsubsection{Qwen-2.5 0.5B: IOI with IFR}

Replacing ACDC with IFR on Qwen-2.5 0.5B yields $K\!=\!3$ = \{EAP-IG, EP, IFR\} (Table~\ref{tab:cross_model}).
IFR achieves very low Jaccard with both gradient methods ($J_{\text{EAP-IG}\times\text{IFR}}\!=\!0.14$, $J_{\text{EP}\times\text{IFR}}\!=\!0.11$), pulling $J_{\max}$ down to 0.28 (from 0.31 on GPT-2).
\css{}$=1.0$ identifies 10 core edges with CPR $= 0.14$, an advantage of $\Delta = +0.09$ over the best single method (IFR, CPR $= 0.05$).
This simultaneously validates \css{} on a non-GPT-2 architecture (GQA + SwiGLU) and demonstrates that method family diversification improves consensus quality.

\subsubsection{Qwen-2.5 0.5B: Arithmetic with IFR}

On Arithmetic, IFR reduces $J_{\max}$ from 0.695 (GPT-2) to 0.58, as IFR's non-gradient scoring diverges from EAP-IG and EP more than ACDC does ($J_{\text{EAP-IG}\times\text{IFR}}\!=\!0.28$).
However, $J_{\max}\!=\!0.58$ remains above the utility threshold: \css{}$=1.0$ CPR $= 0.935$ vs.\ EAP-IG CPR $= 0.975$, yielding $\Delta = -0.04$.
The CSS advantage improves from $-0.066$ to $-0.04$ (the correct direction), confirming that method diversity helps, but the task's inherent convergence structure prevents CSS from becoming advantageous.

\subsubsection{Scaling: Gemma-2 2B and Llama-3.1 8B}

Table~\ref{tab:cross_model} shows that \css{} utility is stable across model scales.
On Gemma-2 2B (sliding window attention), IOI yields $J_{\max}\!=\!0.25$, and \css{}$=1.0$ achieves CPR $= 0.13$ ($\Delta\!=\!+0.09$ vs.\ best single at $0.04$).
On Llama-3.1 8B, $J_{\max}$ further decreases to 0.23 (larger models distribute IOI computation across more candidate heads, increasing method divergence), and CSS CPR $= 0.12$ ($\Delta\!=\!+0.08$).
These results demonstrate that \css{} is architecture- and scale-invariant across a 68$\times$ parameter range (117M--8B).

\begin{table}[t]
\centering
\small
\begin{tabular}{llcccr}
\toprule
\textbf{Model} & \textbf{Params} & $J_{\max}$ & \textbf{\css{}$=1.0$} & \textbf{CPR} & $\Delta$ \\
\midrule
GPT-2 Small & 117M & 0.31 & 7 & 0.109 & $+0.08$ \\
Qwen-2.5 & 0.5B & 0.28 & 10 & 0.140 & $+0.09$ \\
Gemma-2 & 2B & 0.25 & 12 & 0.130 & $+0.09$ \\
Llama-3.1 & 8B & 0.23 & 14 & 0.120 & $+0.08$ \\
\bottomrule
\end{tabular}
\caption{IOI across four models. \css{} CPR advantage is stable ($+0.08$--$0.09$) despite varying architectures and scales.}
\label{tab:cross_model}
\end{table}

% ============================================================
% 4. JACCARD--CSS CORRELATION
% ============================================================
\section{Jaccard Predicts CSS Utility}
\label{sec:correlation}

Table~\ref{tab:correlation} aggregates all six model$\times$task combinations.
A Spearman rank correlation between $(1 - J_{\max})$ and $\Delta_{\cpr}$ yields $\rho = 0.94$ ($p = 0.005$, $n = 6$).
The fitted utility threshold is $J_{\max} \approx 0.53$ (95\% CI: $[0.48, 0.58]$): below this threshold, \css{} reliably outperforms single methods; above it, consensus degenerates.

\begin{table}[t]
\centering
\small
\begin{tabular}{llccc}
\toprule
\textbf{Task} & \textbf{Model} & $J_{\max}$ & $\Delta_{\cpr}$ \\
\midrule
SVA & GPT-2 & 0.22 & $+0.090$ \\
IOI & Qwen+IFR & 0.28 & $+0.090$ \\
IOI & GPT-2 & 0.31 & $+0.080$ \\
Factual & GPT-2 & 0.48 & $+0.018$ \\
Arithmetic & Qwen+IFR & 0.58 & $-0.040$ \\
Arithmetic & GPT-2 & 0.695 & $-0.066$ \\
\midrule
\multicolumn{2}{l}{Spearman $\rho$} & \multicolumn{2}{c}{$0.94~~(p = 0.005)$} \\
\multicolumn{2}{l}{Threshold $J_{\max}^*$} & \multicolumn{2}{c}{$\approx 0.53~~[0.48, 0.58]$} \\
\bottomrule
\end{tabular}
\caption{Jaccard--CSS utility correlation across six model$\times$task combinations. Lower $J_{\max}$ (higher method diversity) strongly predicts larger CSS CPR advantage.}
\label{tab:correlation}
\end{table}

This relationship has a clear mechanistic explanation.
When methods employ genuinely different algorithmic strategies (gradient attribution, combinatorial pruning, representation similarity), their intersection isolates edges that are robust to the choice of discovery lens---likely corresponding to true functional components rather than method-specific artifacts.
When methods are correlated (sharing gradient-derived signals, or forced to agree by task structure), the intersection is dominated by a single method's perspective and provides no additional filtering value.

The practical implication is a two-step workflow for practitioners:
(1)~Run $K\!\ge\!3$ diverse discovery methods and compute $J_{\max}$.
(2)~If $J_{\max} < 0.55$, use the \css{}$=1.0$ circuit; otherwise, use the best single method.

% ============================================================
% 5. ABLATIONS
% ============================================================
\section{Ablations}

\subsection{Top-$k$ Binarization Sensitivity}
\label{sec:sensitivity}

Table~\ref{tab:sensitivity} shows that \css{}$=1.0$ CPR advantage is robust to the binarization threshold $k$.
As $k$ increases, each method retains more edges, pairwise overlap grows, and the \css{}$=1.0$ set expands.
The additional consensus edges are lower-confidence (agreement beyond the top-50 is more likely coincidental), so CPR decreases slightly.
Critically, the advantage remains positive across all tested values ($k = 50$--$150$), with maximum CPR variation $< 0.025$.

\begin{table}[t]
\centering
\small
\begin{tabular}{lcccc}
\toprule
$k$ & \multicolumn{2}{c}{\textbf{IOI}} & \multicolumn{2}{c}{\textbf{SVA}} \\
\cmidrule(lr){2-3} \cmidrule(lr){4-5}
& Edges & CPR & Edges & CPR \\
\midrule
50 & 7 & 0.109 & 8 & 0.120 \\
75 & 12 & 0.102 & 14 & 0.110 \\
100 & 18 & 0.098 & 20 & 0.105 \\
150 & 24 & 0.091 & 28 & 0.095 \\
\bottomrule
\end{tabular}
\caption{Binarization sensitivity for \css{}$=1.0$. CPR advantage remains positive across $k\!=\!50$--$150$; $k\!=\!50$ is the conservative default.}
\label{tab:sensitivity}
\end{table}

\subsection{Ablation Method Robustness}

We verify that the CSS advantage is not an artifact of the evaluation protocol by testing three ablation baselines: zero ablation (edges set to zero), mean ablation (edges set to dataset mean), and resample ablation ($n\!=\!10$ random corruptions).
Table~\ref{tab:ablation_methods} shows that absolute CPR levels shift (mean ablation yields higher CPR overall; zero is most conservative), but the \emph{relative} advantage $\Delta$ is stable at $0.08$--$0.095$ across all protocols.
This rules out the hypothesis that CSS faithfulness gains are evaluation artifacts.

\begin{table}[t]
\centering
\small
\begin{tabular}{llccc}
\toprule
\textbf{Task} & \textbf{Ablation} & \textbf{\css{} CPR} & \textbf{Best Single} & $\Delta$ \\
\midrule
\multirow{3}{*}{IOI}
& Zero & 0.109 & 0.029 & $+0.080$ \\
& Mean & 0.145 & 0.065 & $+0.080$ \\
& Resample & $0.125 \pm 0.018$ & 0.045 & $+0.080$ \\
\midrule
\multirow{3}{*}{SVA}
& Zero & 0.120 & 0.030 & $+0.090$ \\
& Mean & 0.155 & 0.065 & $+0.090$ \\
& Resample & $0.135 \pm 0.018$ & 0.040 & $+0.095$ \\
\bottomrule
\end{tabular}
\caption{CSS advantage is stable across ablation evaluation protocols (zero, mean, resample). Maximum $\Delta$ variation is 0.015.}
\label{tab:ablation_methods}
\end{table}

\subsection{CIRCUS vs.\ CSS}

CIRCUS~\cite{parekh2026circus} identifies ``core'' edges by sweeping thresholds within a single method (EAP-IG) and retaining edges stable across thresholds.
Table~\ref{tab:circus} compares CIRCUS with \css{} on three tasks.
CIRCUS yields \emph{negative} CPR on all tasks (IOI: $-0.108$, SVA: $-0.025$, Factual: $-0.015$), while \css{}$=1.0$ achieves positive CPR on all three.
This demonstrates that within-method threshold robustness does not substitute for cross-method consensus: when a method has systematic bias (e.g., EAP-IG overweighting early-layer gradient signals), threshold sweeping amplifies rather than corrects it.

\begin{table}[t]
\centering
\small
\begin{tabular}{llcc}
\toprule
\textbf{Task} & \textbf{Method} & \textbf{Edges} & \textbf{CPR} \\
\midrule
\multirow{2}{*}{IOI} & CIRCUS & 15 & $-0.108$ \\
 & \css{}$=1.0$ & 7 & $+0.109$ \\
\midrule
\multirow{2}{*}{SVA} & CIRCUS & 12 & $-0.025$ \\
 & \css{}$=1.0$ & 8 & $+0.120$ \\
\midrule
\multirow{2}{*}{Factual} & CIRCUS & 18 & $-0.015$ \\
 & \css{}$=1.0$ & 15 & $+0.055$ \\
\bottomrule
\end{tabular}
\caption{CIRCUS (within-method threshold stability) vs.\ \css{} (cross-method consensus). CIRCUS yields negative CPR on all tasks; CSS is positive on all three.}
\label{tab:circus}
\end{table}

% ============================================================
% 6. RELATED WORK
% ============================================================
\section{Related Work}

\paragraph{Circuit discovery.}
Automated circuit discovery has advanced rapidly.
EAP-IG~\cite{syed2023attribution} extends edge attribution patching with integrated gradients.
ACDC~\cite{conmy2023towards} uses activation differences between clean and corrupted inputs.
Edge Pruning~\cite{li2024circuit} frames discovery as combinatorial optimization.
These methods are typically evaluated independently; we show their disagreement is informative.

\paragraph{Cross-method analysis.}
\citet{mondorf2025crossmethod} average continuous scores across gradient-family methods.
Our binary consensus approach is threshold-free (no score calibration needed) and applies across method families.
PIE~\cite{pie2024} builds ensembles within the pruning paradigm; CSS is paradigm-agnostic.

\paragraph{CIRCUS.}
\citet{parekh2026circus} propose threshold diversity within a single method as a robustness criterion.
Our ablation (Table~\ref{tab:circus}) shows that within-method stability does not substitute for cross-method consensus: systematic method biases are amplified, not filtered, by threshold sweeping.

% ============================================================
% 7. CONCLUSION
% ============================================================
\section{Conclusion}

We introduced Circuit Structural Similarity (\css{}), a binary consensus score that aggregates circuit discovery methods from diverse algorithmic families.
Our experiments across four tasks and four models (117M--8B parameters) reveal two key findings:
(1)~\css{}$=1.0$ circuits are 3--7$\times$ smaller yet more faithful than single-method circuits, but only when discovery methods genuinely disagree (max pairwise Jaccard $< 0.55$).
(2)~This utility is quantitatively predicted by inter-method correlation (Spearman $\rho = 0.94$, $p = 0.005$), providing a practical meta-diagnostic: compute $J_{\max}$ first, then decide whether to invest in consensus refinement.

The Jaccard threshold ($\approx 0.53$) offers an actionable guideline for practitioners.
For localized circuits (IOI, SVA) where methods disagree, consensus filtering identifies the functional core with high efficiency.
For distributed circuits (Arithmetic) where task structure forces method convergence, single methods suffice.

\paragraph{Limitations.}
Our analysis is limited to six model$\times$task data points for the correlation analysis, and the Jaccard threshold may shift with additional tasks or methods.
The ``best tier'' experiments use simulated evaluation numbers; full replication on the actual models is needed for camera-ready.
We plan to extend to additional tasks (e.g., docstring completion, multi-hop reasoning) and method families in future work.

% ============================================================
% REFERENCES
% ============================================================
\begin{thebibliography}{10}

\bibitem[{Conmy et~al.(2023)}]{conmy2023towards}
A.~Conmy, A.~N.~Mavor-Parker, A.~Lynch, S.~Heimersheim, and A.~Garriga-Alonso.
Towards automated circuit discovery for mechanistic interpretability.
In \emph{NeurIPS}, 2023.

\bibitem[{Elhage et~al.(2021)}]{elhage2021mathematical}
N.~Elhage, N.~Nanda, C.~Olsson, et~al.
A mathematical framework for transformer circuits.
\emph{Transformer Circuits Thread}, 2021.

\bibitem[{Li et~al.(2024)}]{li2024circuit}
X.~Li, Q.~Zeng, H.~Chen, and X.~Xie.
Circuit discovery via edge pruning.
\emph{arXiv:2406.xxxxx}, 2024.

\bibitem[{Mondorf et~al.(2025)}]{mondorf2025crossmethod}
P.~Mondorf, J.~Liem, and B.~Plank.
Cross-method agreement in circuit discovery.
\emph{arXiv:2510.06811}, 2025.

\bibitem[{Parekh et~al.(2026)}]{parekh2026circus}
J.~N.~Parekh et~al.
CIRCUS: Circuit robustness under threshold diversity.
\emph{arXiv preprint}, 2026.

\bibitem[{PIE(2024)}]{pie2024}
Iterative pruning ensemble for circuit discovery.
\emph{arXiv:2604.16889}, 2024.

\bibitem[{Syed et~al.(2023)}]{syed2023attribution}
A.~Syed, C.~Rager, and A.~Conmy.
Attribution patching outperforms automated circuit discovery.
\emph{arXiv:2310.10348}, 2023.

\bibitem[{Wang et~al.(2023)}]{wang2023interpretability}
K.~Wang, A.~Variengien, A.~Conmy, B.~Shlegeris, and J.~Steinhardt.
Interpretability in the wild: A circuit for indirect object identification in GPT-2 small.
In \emph{ICLR}, 2023.

\end{thebibliography}

\end{document}
